\documentclass[
  oneside,
  fullwidthstop=false,
  fontset=fandol,
  times=false,
  minted=false, % 改为 true 可启用 Pygments 高亮（需 -shell-escape）
]{tongjithesis}

\usepackage{tikz}
\usetikzlibrary{shapes,arrows,positioning,calc,shadows,trees,arrows.meta,backgrounds,fit}

%%%%%%%%%%%%%%%%%%%%%%%%%%%%%%%%%%%%%%%%%%%%%%%%%%%%%%%%
% 参考文献：默认使用 A。若改用 B，请注释掉 A 的全部命令。
% Citation: method A is active. Comment it out before enabling B.
%%%%%%%%%%%%%%%%%%%%%%%%%%%%%%%%%%%%%%%%%%%%%%%%%%%%%%%%

%%%%%%%%%%%%%%%%%%%%%%%%%%%%%%%%%%%%%%%%%%%%%%%%%%%%%%%%
% A. biblatex + biber（推荐）
%%%%%%%%%%%%%%%%%%%%%%%%%%%%%%%%%%%%%%%%%%%%%%%%%%%%%%%%
\usepackage[
  backend=biber,
  style=gb7714-2025,  % 如需 2015 版，改为 gb7714-2015
  gbnamefmt=lowercase,
  gbpub=false,
  doi=false,
  natbib=true,
]{biblatex}
\addbibresource{bib/note.bib}

%%%%%%%%%%%%%%%%%%%%%%%%%%%%%%%%%%%%%%%%%%%%%%%%%%%%%%%%
% B. BibTeX + gbt7714（传统方式，不要与 A 同时启用）
%%%%%%%%%%%%%%%%%%%%%%%%%%%%%%%%%%%%%%%%%%%%%%%%%%%%%%%%
% \usepackage[sort&compress]{gbt7714}
% \bibliographystyle{gbt7714-numerical}
% \setlength{\bibsep}{3pt}

\begin{document}

% 通用封面配置。留空的可选字段（副标题、系所、学期、班级）不会显示。
% report-type 可按需填写：课程论文、实验报告、课程设计报告、
% 调研报告、读书报告、实习报告、田野调查报告等。
\tongjisetup{
  report-type = {同济大学课程报告},
  title       = {报告题目},
  subtitle    = {},
  title-en    = {Report Title},
  subtitle-en = {},
  school      = {学院或学部},
  department  = {},
  major       = {专业名称},
  course      = {课程名称},
  semester    = {},
  class       = {},
  student-id  = {学号},
  author      = {姓名},
  instructor  = {课程教师},
  date        = {\today},
}

\MakeCover

\clearpage

\pagestyle{firststyle}
% 摘要应能脱离正文，概括背景、对象、方法、结果和结论，通常不超过一页。
% 下方文字与后文示例一致，写稿时请整段替换。
\MakeAbstract{
本报告演示通用课程报告的写法与版式。示例任务是在统一标准下比较三种可替换方案，标准包括完整性、准确性与表达，综合得分由等权加权得到。正文用三线表列出分项结果，用流程图说明工作顺序，并用并排插图对照评价维度与综合得分。示例中的名称和分数只服务排版，正式提交前应换成课程要求的材料、数据与结论。
}{课程报告，三线表，插图，综合评价}


\clearpage
\tableofcontents   % 放置目录
\cleardoublepage

\pagestyle{mainstyle}
\section{概述}
\label{sec:intro}

\subsection{研究背景}
请在此介绍课程背景、选题来源、现实或理论语境，以及本报告关注的核心对象。理工医科可说明问题与应用场景，人文社科可交代议题、材料和研究语境，设计艺术类可说明创作命题与使用场景。

\subsection{问题与目标}
请明确研究问题、实践任务或创作目标，并说明报告的范围与边界。

\subsection{报告结构}
请在此简要介绍本文各章节安排。

% 写作提示：综述与题目直接相关的文献、理论或标准，并说明材料从何而来。
\section{文献、理论与材料}
\label{sec:context}

\subsection{相关研究或理论基础}
不同课程会把本节写成文献综述、理论框架或基础知识。示例保留三条可以并列的路径，并用\cite{gbt7714-2015,lamport1978,vaswani2017,knuth1986,hushidong2025}演示国标、期刊、会议、图书和网络文献的引用。正式写作时，改成与题目直接相关的文献、标准或案例。

\begin{table}[htbp]
  \centering
  \caption{常见分析路径及其适用情形}
  \label{tab:paths}
  \begin{tabular}{@{}l >{\raggedright\arraybackslash}p{0.36\textwidth} >{\raggedright\arraybackslash}p{0.40\textwidth}@{}}
    \toprule
    \heiti 路径 & \heiti 主要做法 & \heiti 更适用的任务 \\
    \midrule
    实验测量 & 控制条件，并记录可以重复的观测 & 对象能够测量，结果需要对照 \\
    文本细读 & 围绕材料做分层解释 & 论点依赖原文、档案或作品 \\
    案例比较 & 用同一标准并列多个对象 & 需要解释差异及其条件 \\
    \bottomrule
  \end{tabular}
\end{table}

\tabref{tab:paths}里的路径可以改成理论流派、实验方案、场地样本或设计方向。表格的作用是把差异放到同一组标准下。

\subsection{研究对象或资料来源}
示例没有使用外部数据。方案 A、B、C 的分项分数写在附录\tabref{tab:raw}，正文只引用汇总结果。若课程要求说明样本、文本、档案或实验材料，应在这里交代来源、筛选条件，以及不能使用的部分。

\subsection{概念界定与评价标准}
三项指标都按 0 到 1 记录，越高越好。完整性表示对象是否被覆盖，准确性表示判断是否有材料支持，表达表示结果是否清楚、可否复核。三项等权，计算方式见\eqref{eq:score}。

% 写作提示：说明总体路线、关键步骤，以及方法为何可靠。
\section{研究或实践方法}
\label{sec:method}

\subsection{总体思路}
先固定评价标准，再按同一流程处理每个方案，最后汇总比较。\figref{fig:flow}是这一顺序的示意图。实验、调查、文本分析和方案设计都可以沿用它，只需替换每一步的具体工作。

\begin{figure}[htbp]
  \centering
  \begin{tikzpicture}[node distance=8mm]
    \node[tongji node] (q) {界定问题};
    \node[tongji node, right=of q] (m) {选择方法};
    \node[tongji node, right=of m] (p) {记录过程};
    \node[tongji node, right=of p] (r) {分析结果};
    \draw[tongji arrow] (q) -- (m);
    \draw[tongji arrow] (m) -- (p);
    \draw[tongji arrow] (p) -- (r);
  \end{tikzpicture}
  \caption{课程报告的基本工作流程}
  \label{fig:flow}
\end{figure}

照片、扫描件或软件导出的图像放在 \texttt{figures/}，并按版心宽度插入，例如：
\begin{center}
  \verb|\includegraphics[width=0.8\linewidth]{figures/example.pdf}|
\end{center}
流程和统计图优先使用矢量图，打印后线条更清楚。

\subsection{方法与步骤}
实施时按下面的顺序推进：
\begin{enumerate}
  \item 写清问题边界和评价指标。
  \item 用同一流程处理每个方案。
  \item 汇总得分，并让公式、表格和插图使用同一组数字。
\end{enumerate}

综合得分把三项分数合成一个结果，便于排序，也保留回到分项的可能：
\begin{equation}
  S = \sum_{i=1}^{n} w_i s_i, \qquad \sum_{i=1}^{n} w_i = 1
  \label{eq:score}
\end{equation}
本例取 $n = 3$，且 $w_i = 1/3$。$s_i$ 是第 $i$ 项得分。计算时保留两位小数，与\tabref{tab:scores}一致。

\begin{thm}
若 $w_i = 1/n$，则综合得分等于各项得分的算术平均。
\end{thm}
\begin{pf}
由\eqref{eq:score}，
\begin{align}
  S &= \sum_{i=1}^{n} \frac{1}{n} s_i \label{eq:mean}\\
    &= \frac{1}{n}\sum_{i=1}^{n} s_i. \nonumber
\end{align}
\end{pf}

\subsection{可靠性与合规性}
示例不涉及个人隐私、实验安全和未公开数据。若实际任务包含问卷、访谈、人体或场地信息，应在这里说明知情同意、匿名化，以及可复现材料放在何处。更换权重或指标时，公式、表格和插图要一起修改\footnote{只改其中一处，正文和附录中的数字就会对不上。}。

% 写作提示：按逻辑或时间顺序写过程，并保留关键异常的处理依据。
\section{分析、实验或实践过程}
\label{sec:process}

\subsection{实施条件}
示例只依赖已经记录的分项分数，不依赖特定软件、仪器或场地。需要重复计算时，可以把步骤写成算法或代码；课程不要求代码时，删去对应环境即可。

\subsection{核心过程}
\algoref{alg:score}是\eqref{eq:score}的逐步形式。算法标题放在上方，输入和输出单独列出。

\begin{algorithm}[htbp]
  \caption{由分项得分计算综合得分}
  \label{alg:score}
  \begin{algorithmic}[1]
    \Require 分项得分 $s_1,\ldots,s_n$ 与权重 $w_1,\ldots,w_n$
    \Ensure 综合得分 $S$
    \State $S \gets 0$
    \For{$i \gets 1,\ldots,n$}
      \State $S \gets S + w_i s_i$
    \EndFor
    \State \Return $S$
  \end{algorithmic}
\end{algorithm}

同一计算也可以写成代码，便于附上能够运行的实现。代码标题放在下方，见\listingref{lst:score}。

\begin{lstlisting}[language=Python,caption={按权重汇总分项得分},label={lst:score}]
def weighted_score(scores, weights):
    total = 0.0
    for score, weight in zip(scores, weights):
        total += score * weight
    return total
\end{lstlisting}

\subsection{关键问题及处理}
示例里没有缺失值。若某一项没有记录，应先决定是补计、剔除，还是单独报告，并在\tabref{tab:scores}的表注里写明。

\section{结果与分析}
\label{sec:results}

\subsection{主要结果}
请客观呈现研究发现、实验数据、系统表现、调查结果、论证结论或作品成果。

\subsection{结果分析}
请结合前文方法与评价标准解释结果，可使用统计分析、案例比较、理论阐释、误差分析或设计评估。

\subsection{验证或论证}
请说明结果如何回应研究问题或课程任务，并提供必要的对照、证据、引用或验证过程。

% 写作提示：把结果放回文献、任务目标和适用范围里讨论，并写明局限。
\section{讨论}
\label{sec:discussion}

\subsection{结果讨论}
方案 C 的综合得分最高，主要因为完整性和表达两项领先；方案 B 在准确性上更好。若课程更重视可以核对的依据，就不该只按\tabref{tab:scores}的最后一列排序，而应提高准确性的权重，并同步修改\eqref{eq:score}和\figref{fig:results}。

\subsection{局限与改进}
示例只有三个方案和三项指标，没有误差、重复测量或外部对照。分数是事先设定的，不能当成实验发现。正式报告应说明样本范围和材料缺口，并估计增加指标或改变权重之后，排序可能如何变化。

\begin{rem}
同一指标在正文、表格和插图中应使用相同的名称和精度。
\end{rem}

% 写作提示：回答最初的问题，归纳主要工作，并指出后续方向。结论中不引入正文没有讨论过的新材料。
\section{总结}
\label{sec:conclusion}

示例任务在等权标准下把方案 C 排在前面，同时保留了方案 B 准确性更高这一分项差异。公式、三线表、流程图和分组插图使用的是同一组数字，因此可以互相核对。换成真实题目时，应保留这种前后对应，并删去与课程无关的示例。


%%%%%%%%%%%%%%%%%%%%%%%%%%%%%%%%%%%%%%%%%%%%%%%%%%%%%%%%
% 参考文献列表生成（选择与上方相同的选项A或B）
% Reference list generation (choose the same option A or B as above)
%%%%%%%%%%%%%%%%%%%%%%%%%%%%%%%%%%%%%%%%%%%%%%%%%%%%%%%%

% A. biblatex
\clearpage
\printbibliography[heading=bibintoc,title={参考文献}]

% B. BibTeX。启用前请先注释掉上面的 \printbibliography。
% \addcontentsline{toc}{section}{参考文献}
% \bibliography{bib/note.bib}


\clearpage
% 写作提示：放置推导、问卷、访谈提纲、原始数据、图版、代码或配置。正文用 \tabref 等方式引用。
\appendix
\section{附录}
\label{sec:appendix}

\subsection{补充材料}
\tabref{tab:scores}中的综合得分来自下面的原始记录。问卷、推导、图版或更长的数据也可以放在本节，并在正文中引用。

\begin{table}[htbp]
  \centering
  \caption{分项得分的原始记录}
  \label{tab:raw}
  \begin{tabular}{lccc}
    \toprule
    方案 & 完整性 & 准确性 & 表达 \\
    \midrule
    方案 A & 0.72 & 0.68 & 0.75 \\
    方案 B & 0.81 & 0.84 & 0.79 \\
    方案 C & 0.86 & 0.83 & 0.88 \\
    \bottomrule
  \end{tabular}
\end{table}

\end{document}
